% intro.tex — Introduction.
% Empirical RViT+ extension of Morgan, Albanna & Herman (2025), and companion
% to the separate 2026 normative VDA manuscript. THE model is affine_ew.

\section{Introduction}
\label{sec:intro}

For a quarter of a century, a single question has organised the study of
visual attention: \emph{how} does attention improve perception? Attending a
location makes an observer faster and more accurate at the events that occur
there, and the size of that benefit has been mapped in exhaustive
psychophysical and physiological detail \cite{carrasco2011_visual_attention_25y,
posner1980_orienting}. What the benefit \emph{is}---the computation that turns a
spatial cue into better vision---has proven far harder to pin down. Reviewing
the field, Carrasco \cite{carrasco2011_visual_attention_25y} frames three
standing accounts. On the first, attention \emph{enhances the signal}: it
multiplies the gain of the neurons that encode the attended stimulus, sharpening
their response the way a contrast increase would
\cite{reynolds_heeger2009_normalization, mcadams_maunsell1999_v4_tuning,
treue_martinez_trujillo1999_feature_attention, sani2017_temporal_v4_gain}. On
the second, attention improves sensitivity by \emph{reducing external noise}, or
equivalently by selecting the relevant input more efficiently and excluding
distractors and competing locations \cite{CohenMaunsell2009}. On the third, the
gain in performance is not perceptual at all but sits at the \emph{decision
stage}---a shift in the observer's response criterion, or a reduction in spatial
uncertainty about where the relevant event will occur \cite{MullerFindlay1987,
posner1980_orienting}.

These are not idle alternatives. They make different predictions about which
neurons should change, about how the psychometric and the criterion should move,
and about where in the processing stream an intervention should have its effect.
Yet Carrasco's own resolution is appropriately plural: the three accounts,
she writes, ``are not mutually exclusive; rather, all three are likely to
contribute'' \cite{carrasco2011_visual_attention_25y}. Behaviour is a summary of a
system with many internal stages, and a change in hit rate can be consistent with
a gain change, a noise change, or a criterion change. Signal-detection analysis can,
with care, separate a sensitivity ($\dprime$) effect from a criterion ($c$) effect
\cite{MullerFindlay1987}; dedicated external-noise paradigms are needed to adjudicate
signal enhancement against external-noise reduction. An instrumented model adds a
different form of access: its internal representation can be read out and its
attention can be injected or suppressed at a chosen location. That access supports
direct tests of criterion, sensitivity, and spatial gating, but it does not by itself
turn those measurements into a quantitative decomposition of all three accounts.

A computational model in which attention is learned rather than prescribed offers
that access. Its internal states can be decoded and its attention machinery clamped
or injected. The resulting tests can establish whether a manipulation moves
criterion and sensitivity and whether its behavioural effect is spatially gated.
They cannot estimate a three-way contribution fraction without an explicit
generative decomposition and a dedicated external-noise manipulation, neither of
which is part of the present battery.

Here we apply those direct tests. We study a single recurrent vision transformer
equipped with a spatial working-memory state, trained by reinforcement learning
together with a temporal joint-embedding predictive (JEPA) auxiliary objective
(loss weight $0.5$) to
perform a cued change-detection task: a spatial cue
indicates where a change is likely, a brief blank delay follows, an array of
oriented Gabor stimuli appears, and the model is rewarded for reporting a change
in one of them and for withholding on unchanged displays. The model is never
told to attend, never given a hand-specified saliency map, and never trained on
neural or behavioural targets. The JEPA head has no direct forward input to the
decision, but its weighted loss backpropagates through and shapes the recurrent
representation. Under this joint objective the model develops a graded internal
attention map over the stimulus array---a priority map. It orients at the cue,
relaxes near baseline during the blank, and later re-acquires the cued location;
that sequence is compatible with an intervening memory trace, but no cue-location
decoder across the delay is reported. We then apply to this model the apparatus that
the biological literature can only partially bring to bear: signal-detection
read-outs of sensitivity and criterion, graded clamping of the attention map
onto or off any location, causal microstimulation that injects attention where
none was deployed, and time-resolved decoding of task variables from the memory
state. The core causal battery is performed on one canonical affine-feedback $d_{\mathrm{mem}}=128$
checkpoint. Set-size and paradigm comparisons use separately trained checkpoints of the
same architecture; we identify those comparisons explicitly rather than treating them as
repeated measurements of one set of weights.

The central result is that this one model makes criterion, sensitivity, and spatial
gating measurable under direct manipulation. Injecting attention onto an unchanged
location on catch trials evokes false detections preferentially at the monitored
location, a spatially gated effect
(Section~\ref{sec:detection}). Attention is also associated with efficient selection:
a cued-versus-uncued benefit is present in the four-stimulus checkpoint, while the
single-stimulus task has no comparable uncued active item
(Sections~\ref{sec:detection}, \ref{sec:vda-battery}). This is compatible with a
noise-limited detector but is not a pure set-size intervention or an external-noise test. Attention also moves
the decision stage: in the canonical VDA4 sweep, moving the nominal clamp from
suppress to enhance shifts the sampled response criteria toward liberal responding,
the textbook attentional criterion effect; the VDA9 curve is not uniformly monotone
(Section~\ref{sec:detection}). Clamping attention while reading out
signal-detection quantities directly operationalizes both the criterion and
sensitivity responses; and a
misdirection probe reveals a failure mode invisible to behaviour alone, in which
aiming the model's attention at the wrong location makes it \emph{miss} a real
change while becoming \emph{more} confident that nothing happened
(Section~\ref{sec:detection}). This is a model-level account of how attentional
misallocation manufactures confident perceptual errors, and it is the kind of
mechanistic claim that a scalar behavioural benefit can never license.

The same affine-feedback architecture and clamp analysis are then applied to separately
trained paradigm checkpoints. In a preliminary, cued Luo--Maunsell-style design
that separates a reward regime rewarding sensitivity from one rewarding a
criterion shift \cite{luo_maunsell2018_criterion_sensitivity}, the reward-regime comparison shows the expected contrast---the sensitivity regime moves $\dprime$ while the criterion
regime moves bias---and the follow-up localises the criterion effect to the
decision stage, a shift in the actor's declare-bias, while the sensitivity
effect is carried by attentional gain
(Section~\ref{sec:luomaunsell}). In a separately trained checkpoint where the cue carries spatial validity rather
than value, the architecture \emph{uses} that validity gradedly, in the Posner manner:
its cueing benefit grows with cue reliability and it diverts resources away from
the uncued location as the cue becomes more predictive
(Section~\ref{sec:value}). And where the primate literature reports that value
sensitivity is absolute rather than relative to context
\cite{baruni2015}, the model agrees, yielding a principled null
that matches the regime the companion normative analysis predicts
(Section~\ref{sec:value}).

Beyond the mechanisms of the detection benefit itself, the model's
working memory lets us engage a second, tightly related literature: the claim
that spatial attention and spatial working memory interact, with attention
serving as the rehearsal mechanism that maintains a remembered location across a
delay \cite{awh_jonides2001_overlapping_attention_wm, awh1998_rehearsal,
kane2001_controlled_attention}. In the model, value is latched at the cue and held across the delay, while change
presence and location become decodable at the observed change frame
(Section~\ref{sec:detection})---a candidate functional analogue of an event-location
signal in posterior parietal cortex
\cite{bisley_goldberg2010}. Because the memory has an
explicit capacity---the width of its state---we can ask whether capacity governs
controlled, cue-appropriate detection the way working-memory span predicts
controlled attention in humans (Section~\ref{sec:capacity}), and because cued-location activity during the blank is compatible with maintenance, we can ask
whether it is load-bearing rehearsal by probing and lesioning it during the delay
(Section~\ref{sec:rehearsal}). No delay-probe or causal delay-clamp result is reported
here; rehearsal remains a hypothesis.

This paper extends the 2025 empirical Recurrent ViT of Morgan, Albanna, and Herman
\cite{morgan_albanna_herman2025_rvit}. It is also the empirical companion to a
\emph{separate} 2026 normative signal-detection treatment
\cite{morgan_albanna_herman2026_vda}, rather than a revision of that paper. Where the
normative companion asks, from first principles, \emph{when} value-directed
re-allocation of attention should gain over an observer that has already
optimised its decision criteria, the present paper asks what a system that
learns the task \emph{de novo} actually does, and whether it recapitulates the
signatures the biological literature has catalogued. The two meet at a shared conceptual question: a cued detection task in which
criterion adjustment and attentional re-allocation can provide distinct routes to
reward, and in which telling them apart requires more than the behaviour they jointly
produce. The normative account states conditions under which the distinction should
matter; the empirical checkpoints studied here make both routes measurable. Agreement
between the papers is a consistency relation, not a claim that the empirical network
implements the normative derivation exactly.

The remainder of the paper is organised by the question each group of
experiments investigates. Section~\ref{sec:model} describes the model and the
cued change-detection environment. Section~\ref{sec:detection} presents the core
detection results---the psychometric and reaction-time signatures, the learned
priority map and its time course, the signal-detection dissociation of criterion
from sensitivity, the misdirection failure mode, the time-resolved decoding of
task variables from memory, and the causal microstimulation of the priority
map---establishing directly measured criterion and sensitivity effects plus spatial
gating, while leaving external-noise reduction untested. Section~\ref{sec:value} takes up value and validity: how the
model uses a purely spatial validity cue, and why it treats value as absolute
rather than relative. Section~\ref{sec:luomaunsell} reproduces the
sensitivity/criterion double dissociation in a canonical primate design and
localises each effect to its processing stage. Sections~\ref{sec:vda-battery},
\ref{sec:capacity}, and~\ref{sec:rehearsal} then use the model's working memory
to engage the noise-limited set-size account, the working-memory-capacity view
of controlled attention, and the attention-as-rehearsal hypothesis in turn. The
Discussion returns to Carrasco's framework and distinguishes what the present
battery operationalizes from what still requires a dedicated external-noise design.
